\section{Cách tiếp cận đa tín hiệu}
Một method đơn lẻ thường không đủ kết luận chắc chắn. Ví dụ, ảnh thật được nén lại nhiều lần có thể mất metadata; ảnh AI được chỉnh sửa hậu kỳ có thể có artifact giống ảnh thật; ảnh chụp màn hình có thể làm thay đổi dấu vết nén. Vì vậy, SourceVerify áp dụng cách tiếp cận đa tín hiệu: nhiều method cùng phân tích một tệp, sau đó hệ thống tổng hợp điểm.

Ưu điểm của hướng đa tín hiệu là:
\begin{itemize}
  \item Giảm phụ thuộc vào một dấu hiệu duy nhất.
  \item Dễ giải thích kết quả cho người dùng.
  \item Có thể mở rộng dần bằng cách thêm method mới.
  \item Phù hợp với giai đoạn Project I vì không bắt buộc phải có tập dữ liệu huấn luyện lớn.
\end{itemize}
